\mychapter{VACIAGO}
\section{Fondations}

Digital evidence refers to information stored or transmitted in digital form that may have evidential value within an investigation or judicial proceeding. Several definitions have been proposed

\medskip

\noindent
A commonly used definition provided by the \textit{Scientific Working Group on Digital Evidence (SWGDE)} states that digital evidence is:

\begin{quote}
\small
``Any information of evidential value stored or transmitted in digital form.''
\end{quote}

\medskip

\noindent
Similarly, Eoghan Casey defines digital evidence as:

\begin{quote}
\small
``Any probative information stored or transmitted in digital form that a party to a court case may use at trial.''
\end{quote}

\medskip

\noindent
Another institutional definition provided by the \textit{Council of Europe} describes electronic evidence as information generated, stored or transmitted through electronic devices that may be relied upon in judicial proceedings.
\bigskip

\noindent
\textbf{Digital forensics} can be defined as the discipline concerned with the identification, acquisition, preservation, and analysis of digital evidence in a manner that maintains its integrity and ensures its admissibility in legal proceedings.

\subsection{Digital Forensics vs Computer Forensics}

Computer forensics traditionally focuses on evidence originating from computers, whereas digital forensics includes evidence produced by a wide range of digital systems.

\medskip

\noindent
Sources of digital evidence today include:

\begin{itemize}
\item personal computers
\item smartphones and tablets
\item cloud services
\item network infrastructures
\item Internet platforms and social media
\item Internet of Things (IoT) devices
\end{itemize}
\bigskip

\noindent
Digital evidence presents several intrinsic characteristics that distinguish it from traditional physical evidence and require specialized forensic methodologies.

\begin{itemize}
\item \textbf{Invisibility}: digital evidence is not directly observable and requires technical tools for analysis.
\item \textbf{Expert interpretation}: its meaning often requires interpretation by digital forensic specialists.
\item \textbf{Fragility}: digital data may be altered or destroyed through normal system operations.
\item \textbf{Replicability}: digital information can be copied indefinitely without degradation.
\end{itemize}

\medskip

\noindent
These characteristics explain why strict technical and procedural safeguards are required during digital investigations.

\medskip

\subsubsection{Legal Requirements of Digital Evidence}

\noindent
For digital evidence to be used within legal proceedings, several legal requirements must be satisfied.

\medskip

\begin{itemize}
\item \textbf{Admissibility}: evidence must be collected in compliance with the law and recognized forensic procedures.
\item \textbf{Authenticity}: the evidence must be protected from alteration or tampering.
\item \textbf{Reliability}: the methods used for acquisition and analysis must ensure trustworthy results.
\item \textbf{Understandability}: the evidence must be presented in a form that can be clearly understood by the judge.
\item \textbf{Proportionality}: investigative measures must respect the fundamental rights of the individuals involved.
\end{itemize}

\medskip

\subsubsection*{Admissibility}

\noindent
Admissibility represents one of the most critical legal concepts in digital forensics. The mere existence of incriminating digital information does not automatically allow its use in court.

\medskip

\noindent
Evidence must be obtained through legally authorized procedures, such as judicial authorization or prosecutorial orders. If investigators acquire evidence without respecting these legal grounds, the evidence may be declared inadmissible even if it clearly proves the commission of a crime.

\medskip

\noindent
This requirement exists not to protect criminals but to safeguard fundamental rights and democratic principles. Investigative powers must always be exercised within the boundaries established by law.

\medskip

\subsubsection*{Authenticity and Reliability}

\noindent
Digital evidence must also guarantee authenticity and reliability. Investigators must ensure that data has not been altered during acquisition or analysis.

\subsubsection*{Understandability}

The ability to translate complex technical analysis into clear explanations is therefore an essential skill in digital forensic practice.

\subsection{The Principle of Proportionality}

\noindent
Digital investigations often involve extensive access to personal information stored within digital devices. For this reason, the principle of proportionality plays a fundamental role.

\medskip

\noindent
This principle requires investigators to limit their analysis strictly to information relevant to the case under investigation.

\medskip

\noindent
During a forensic analysis, investigators may encounter information unrelated to the crime, such as:

\begin{itemize}
\item private communications
\item personal relationships
\item browsing history
\item personal habits
\end{itemize}

\medskip

\noindent
If such information is not relevant to the legal investigation, it should not be disclosed or used in the proceedings. The investigator's task is to determine whether a crime has been committed, not to evaluate the personal life or moral behavior of the individual.

\medskip

\noindent
Failure to respect proportionality may expose investigators to legal liability, including potential civil actions for privacy violations.

\subsection{Artificial Intelligence and Human Oversight}

\noindent
The increasing use of artificial intelligence in digital investigations introduces additional legal and ethical considerations.

\medskip

\noindent
European regulatory frameworks, including the GDPR and the AI Act, emphasize that automated decision-making systems cannot operate without meaningful human oversight.

\medskip

\noindent
Consequently, digital forensic methodologies must balance the efficiency of automated analytical tools with human supervision capable of ensuring accountability, transparency and respect for fundamental rights.




\subsection{Categories of Digital Evidence}

\noindent
Digital evidence can be classified according to the origin of the data and the entity responsible for its creation. A common classification identifies three main categories.

\medskip

\subsubsection{Evidence Created by Humans}

\noindent
The first category includes digital data generated directly through human actions. These data result from intentional interaction with a digital system.

\medskip

\noindent
Two common forms can be distinguished:

\begin{itemize}
\item \textbf{Human-to-human communication}: information exchanged between individuals through digital systems (e.g., e-mail messages or online communications).
\item \textbf{Human-to-computer interaction}: digital artefacts produced when a user interacts with software applications (e.g., creating a text document or editing a file).
\end{itemize}

\medskip

\noindent
In these cases the human actor is directly responsible for the generation of the digital artefact.

\medskip

\subsubsection{Evidence Created Automatically by Computer Systems}

\noindent
A second category includes digital evidence produced automatically by computer systems without direct human intervention.

\medskip

\noindent
These data are generated through automated processes or algorithms executed by software systems. Typical examples include:

\begin{itemize}
\item system log files
\item network traffic logs
\item telephone call records
\item Internet Service Provider records
\end{itemize}


\subsubsection{Evidence Created by Both Humans and Computer Systems}

\noindent
The third category includes digital artefacts generated through the interaction between human input and automated computer processing.

\medskip

\noindent
A typical example is an electronic spreadsheet in which a user enters data while the software automatically processes the information through formulas or computational functions.

\medskip

\noindent
In this case, the resulting digital evidence reflects both the human input and the automated processing performed by the system.

\medskip


\subsection{The ``Big Five'' Principles of Digital Forensics}

\noindent
The Council of Europe has identified a set of fundamental principles (``Big Five'') that guide the correct handling of digital evidence during investigations.

\medskip

\subsubsection*{Data Integrity}

\noindent
No investigative action should modify electronic devices or digital media that may subsequently be presented as evidence in court.

\medskip

\subsubsection*{Chain of Custody}

\noindent
All actions performed on digital evidence must be documented through a traceable audit trail that records how the evidence has been collected, handled, and preserved.

\medskip

\subsubsection*{Specialist Support}

\noindent
When investigations involve complex digital systems, it may be necessary to involve specialized forensic experts capable of conducting technically sound analyses.

\medskip

\subsubsection*{Appropriate Training}

\noindent
Personnel involved in the acquisition and preservation of digital evidence must receive adequate training in forensic procedures in order to avoid accidental alteration of evidence.

\medskip

\subsubsection*{Legality}

\noindent
The authorities responsible for the investigation must ensure that all forensic activities are conducted in accordance with the applicable legal framework and procedural requirements.


\subsection{Standards for Digital Evidence Handling}

\noindent
Digital forensic investigations are increasingly guided by international technical standards that provide methodological guidance for the acquisition, preservation and analysis of digital evidence.

\medskip

\noindent
Among the most relevant frameworks used in practice are standards developed by:

\begin{itemize}
\item the \textbf{National Institute of Standards and Technology (NIST)} in the United States;
\item the \textbf{International Organization for Standardization (ISO)} in the European and international context.
\end{itemize}

\medskip

\noindent
These standards do not constitute binding legal rules but represent internationally recognized technical references that facilitate cooperation and methodological consistency across jurisdictions.

\medskip

\begin{quotation}
  \textbf{Author's note:} \textit{ Professor Vaciago underlines that the goal is not to memorize a specific standard, but to recognize the principal steps involved.}
\end{quotation}


\medskip


\subsubsection{Digital Forensics within Incident Response}

\noindent
Within incident response activities, digital forensic investigations generally follow a structured process composed of four main phases.

\medskip

\begin{itemize}
\item \textbf{Collection:} identification and acquisition of relevant data sources while preserving forensic integrity.
\item \textbf{Examination:}extraction and technical assessment of information from the collected data using validated forensic tools.
\item \textbf{Analysis:} interpretation of the examined data in order to reconstruct events and determine the causes of the incident.
\item \textbf{Reporting:} documentation of the investigative findings, including the procedures followed and the integrity verification of the evidence.
\end{itemize}

\subsubsection{Integrity Assurance}

\noindent
Preserving the integrity of digital evidence is a fundamental requirement during forensic acquisition. Several technical mechanisms are commonly adopted to ensure that evidence remains unaltered throughout the investigative process.

\medskip

\subsubsection*{Bitstream Imaging}

\noindent
Bitstream imaging consists of creating a sector-by-sector copy of the entire storage medium.  
This approach allows investigators to reproduce the exact content of the original device while performing analysis on a duplicate copy.

\medskip

\subsubsection*{Write-Blocking Mechanisms}

\noindent
Write-blocking technologies—either hardware or software—prevent any modification of the original storage medium during the acquisition phase.

\medskip

\noindent
This mechanism ensures that the forensic process does not inadvertently alter the evidence.

\medskip

\subsubsection*{Cryptographic Hash Verification}

\noindent
Cryptographic hashing techniques are used to verify that the copied data is identical to the original evidence.

\medskip

\noindent
Hash functions generate a unique digital fingerprint of the data. By comparing hash values before and after acquisition, investigators can confirm that the integrity of the evidence has been preserved.

\medskip


\subsubsection{Media Sanitisation as a Risk Management Process}

\noindent
Another relevant aspect addressed in forensic and security standards concerns the sanitisation of digital media.

\medskip

\noindent
Media sanitisation refers to procedures designed to prevent the recovery of sensitive information from storage devices after their disposal or reuse.

\medskip

\noindent
Standards such as NIST SP 800-88 describe different sanitisation methods depending on the required level of protection.

\medskip

\begin{itemize}
\item \textbf{Clear:} logical techniques that remove data while allowing the storage medium to remain usable.
\item \textbf{Purge:} advanced sanitisation methods designed to prevent data recovery through sophisticated laboratory techniques.
\item \textbf{Destroy:} physical destruction of the storage medium when absolute data elimination is required.
\end{itemize}

\medskip

\noindent
These procedures are often evaluated within a broader \textit{risk management approach}, where investigators must balance the need for evidence preservation with operational constraints such as time, data volume, and investigative priorities.
